\newif\ifglobal

%%%%%%%%%%%%%%%%%%%%%%%
%%% Compile to view %%%
%%%%%%%%%%%%%%%%%%%%%%%

%%%%%%%%%%%%%%%%%%%%%%%%%%%%%%%%%%%%%%%%%%%%%%%%%%%
%%% Readme:                                     %%%
%%% https://www.overleaf.com/read/bwdjvfjksvjy/ %%%
%%%%%%%%%%%%%%%%%%%%%%%%%%%%%%%%%%%%%%%%%%%%%%%%%%%

% >>> M?: marks a module setting number or important notice

% >>> M4: Set {./relative-path} to external files for this module <<<
\def \modulefiles {./36-TWAI}
% To add an external file, use:
% (Replace '---' with actual filename)
% '\input{\modulefiles/tables/---.tex}' for tables
% '\includegraphics{\modulefiles/figures/---}' for figures
% '\subfile{\modulefiles/---}' for register files


% >>> M1: This module is in global project? <<<
% 'YES' - uncomment; 'NO' - comment out
\globaltrue

\ifglobal
    \documentclass[main\_\_EN.tex]{subfiles}
   
\else
    \documentclass[openany, 10pt]{book}

    % >>> M2: In ./00-shared/preamble-trm-module, also find
    %         settings for visibility of todo notes, watermark <<<
    \usepackage{preamble-trm-module}

    % >>> M3: Temporary labels in English,
    %         you can add your own labels there <<<
    \myexternaldocument{./00-shared/temp-labels-en}

\fi %\ifglobal


\setcounter{secnumdepth}{4}


% >>> M5: If this module needs any updates in
% './00-shared/preamble-trm-module.sty'
% (include additional package, etc.), contact your module PM
% or create the file './00-shared/changelog.tex' make the
% necessary updates yourself and document them in this file <<<


\begin{document}


% Sets language into which pre-defined bits of text are translated
\selectlanguage{english}

% Do not delete this block
\ifglobal
% Add the following front matter if
% this module is outside of global project
\else
    \InputIfFileExists{readme}{}{}
    {\let\clearpage\relax \listoftodos}
    \tableofcontents
    \listoftables
    \listoffigures
\fi


%%%%%%%%%%%%%%%%%%%%%%%%%
%%% TRM Module Begins %%%
%%%%%%%%%%%%%%%%%%%%%%%%%

\hypertarget{twai}{}
\chapter{Two-wire Automotive Interface (TWAI)}
\label{mod:twai}


The Two-wire Automotive Interface (TWAI$^\circledR$) is a multi-master, multi-cast communication protocol with functions such as error detection and signaling and inbuilt message priorities and arbitration. The TWAI protocol is suited for automotive and industrial applications (see Section \ref{sec:twai-functional-protocol} for more details).

\chipname{} contains two TWAI controllers named TWAI 0 and TWAI 1. Each controller can individually be connected to a TWAI bus via an external transceiver. The TWAI controllers contain numerous advanced features, and can be utilized in a wide range of use cases such as automotive products, industrial automation controls, building automation, etc.

\section{Features}
Each of the two TWAI controllers on \chipname{} support the following features:

\begin{itemize}
    \item Compatibility with ISO 11898-1 protocol (CAN Specification 2.0)
    \item Standard Frame Format (11-bit ID) and Extended Frame Format (29-bit ID)
    \item Bit rates from 1 Kbit/s to 1 Mbit/s
    \item Multiple modes of operation:
    \begin{itemize}
        \item Normal
        \item Listen-only (no influence on bus)
        \item Self-test (no acknowledgment required during data transmission)
    \end{itemize}
    \item 64-byte Receive FIFO
    \item Special transmissions:
    \begin{itemize}
        \item Single-shot transmissions (does not automatically re-transmit upon error)
        \item Self Reception (the TWAI controller transmits and receives messages simultaneously)
    \end{itemize}
    \item Acceptance Filter (supports single and dual filter modes)
    \item Error detection and handling:
    \begin{itemize}
        \item Error Counters
        \item Configurable Error Warning Limit
        \item Error Code Capture
        \item Arbitration Lost Capture
    \item Automatic Transceiver Standby
    \end{itemize}
\end{itemize}

\section{Protocol Overview}
\label{sec:twai-functional-protocol}

\subsection{TWAI Properties}

The TWAI protocol connects two or more nodes in a bus network, and allows nodes to exchange messages with deterministic delay. A TWAI bus has the following properties:

{\bfseries Single Channel and Non-Return-to-Zero:} The bus has only one transmission line for single-channel communication, which is half-duplex. Synchronization is also implemented in this channel, so extra channels (e.g., clock or enable) are not required. The bit stream of a TWAI message is encoded using the Non-Return-to-Zero (NRZ) method.

{\bfseries Bit Values:} The single channel can either be in a dominant or recessive state, representing a logical 0 and a logical 1 respectively. A node transmitting data in a dominant state always overrides the other node transmitting data in a recessive state. The physical implementation on the bus is left to the application level to decide (e.g., differential pair or a single wire).

{\bfseries Bit Stuffing:} Certain fields of TWAI messages are bit-stuffed. A transmitter that transmits five consecutive bits of the same value (e.g., dominant value or recessive value) should automatically insert a complementary bit. Likewise, a receiver that receives five consecutive bits of the same value should treat the next bit as a stuffed bit. Bit stuffing is applied to the following fields: SOF, arbitration field, control field, data field, and CRC sequence (see Section \ref{sec:twai-messages} for more details).

{\bfseries Multi-cast:} All nodes receive the same bits as they are connected to the same bus. Data is consistent across all nodes unless there is a bus error (see Section \ref{sec:twai-errors} for more details).

{\bfseries Multi-master:} Any node can initiate a transmission. If a transmission is already ongoing, a node will wait until the current transmission is over before initiating a new transmission.

{\bfseries Message Priority and Arbitration:} If two or more nodes simultaneously initiate a transmission, the TWAI protocol ensures that one node will win arbitration of the bus. The arbitration field of the message transmitted by each node is used to determine which node will win the arbitration.

{\bfseries Error Detection and Signaling: } Each node actively monitors the bus for errors, and signals the detected errors by transmitting an error frame.

{\bfseries Fault Confinement:} Each node maintains a set of error counters that are incremented/decremented according to a set of rules. When the error counters surpass a certain threshold, the node will automatically eliminate itself from the network by switching itself off.

{\bfseries 
Configurable Bit Rate:} The bit rate for a single TWAI bus is configurable. However, all nodes on the same bus must operate at the same bit rate.

{\bfseries Transmitters and Receivers:} At any point in time, a TWAI node can either be a transmitter or a receiver.
\begin{itemize}
    \item A node generating a message is a transmitter. The node remains a transmitter until the bus is idle or until the node loses arbitration. Please note that there could be multiple nodes that act as transmitters during arbitration.
    \item All nodes that are not transmitters are receivers.
\end{itemize}

\subsection{TWAI Messages}
\label{sec:twai-messages}

TWAI nodes use messages to transmit data, and signal errors to other nodes when detecting errors on the bus. Messages are split into various frame types, and some frame types will have different frame formats. 

The TWAI protocol has the following frame types:

\begin{itemize}
    \item Data frame
    \item Remote frame
    \item Error frame
    \item Overload frame
    \item Interframe space
\end{itemize}

The TWAI protocol supports the following frame formats:

\begin{itemize}
    \item Standard Frame Format (SFF) that uses a 11-bit identifier
    \item Extended Frame Format (EFF) that uses a 29-bit identifier
\end{itemize}

\subsubsection{Data Frames and Remote Frames}

Data frames are used by nodes to send data to other nodes, and can have a payload of 0 to 8 data bytes. Remote frames are used for nodes to request a data frame with the same identifier from other nodes, and thus they do not contain any data bytes. However, data frames and remote frames share many fields. Figure \ref{fig:twai-bit-fields} illustrates the fields and sub-fields of different frames and formats.

\begin{figure}[ht!]
    \centering
    \fbox{\includegraphics[width=0.8\textwidth]{\modulefiles/figures/36-TWAI-bit-fields__EN}}
    \caption{Bit Fields in Data Frames and Remote Frames}
    \label{fig:twai-bit-fields}
\end{figure}

{\bfseries Arbitration Field}\\
When two or more nodes transmit a data or remote frame simultaneously, the arbitration field is used to determine which node will win arbitration of the bus. In the arbitration field, if a node transmits a recessive bit while detecting a dominant bit, it indicates that another node has overridden its recessive bit. Therefore, the node transmitting the recessive bit has lost arbitration of the bus and should immediately switch to be a receiver.

The arbitration field primarily consists of a frame identifier that is transmitted from the most significant bit first. Given that a dominant bit represents a logical 0, and a recessive bit represents a logical 1:

\begin{itemize}
    \item A frame with the smallest ID value always wins arbitration.
    \item Given the same ID and format, data frames always prevail over remote frames due to their RTR bits being dominant.
    \item Given the same first 11 bits of ID, a Standard Format Data Frame always prevails over an Extended Format Data Frame due to its SRR bits being recessive.
\end{itemize}

{\bfseries Control Field}\\
The control field primarily consists of the DLC (Data Length Code) which indicates the number of payload data bytes for a data frame, or the number of requested data bytes for a remote frame. The DLC is transmitted from the most significant bit first.

{\bfseries Data Field}\\
The data field contains the actual payload data bytes of a data frame. Remote frames do not contain any data field.

{\bfseries CRC Field}\\
The CRC field primarily consists of a CRC sequence. The CRC sequence is a 15-bit cyclic redundancy code calculated from the de-stuffed contents (everything from the SOF to the end of the data field) of a data or remote frame.

{\bfseries ACK Field}\\
The ACK field primarily consists of an ACK Slot and an ACK Delim. The ACK field indicates that the receiver has received an effective message from the transmitter.

\input{\modulefiles/tables/36-TWAI-data-frames-and-remote-frames__EN}

\subsubsection{Error and Overload Frames}

{\bfseries Error Frames}

Error frames are transmitted when a node detects a bus error. Error frames notably consist of an Error Flag which is made up of six consecutive bits of the same value, thus violating the bit-stuffing rule. Therefore, when a particular node detects a bus error and transmits an error frame, all other nodes will then detect a stuff error and transmit their own error frames in response. This has the effect of propagating the detection of a bus error across all nodes on the bus. 

When a node detects a bus error, it will transmit an error frame starting from the next bit. However, when a node detects a CRC error, the error frame will start at the bit following the ACK Delim (see Section \ref{sec:twai-errors} for more details). The following Figure \ref{fig:twai-fields-of-an-error-frame} shows different fields of an error frame:

\begin{figure}[ht]
    \centering
    \fbox{\includegraphics[width=0.8\textwidth]{\modulefiles/figures/36-TWAI-fields-of-an-error-frame__EN}}
    \caption{Fields of an Error Frame}
    \label{fig:twai-fields-of-an-error-frame}
\end{figure}

\input{\modulefiles/tables/36-TWAI-error-frame__EN}

{\bfseries Overload Frames}\\

An overload frame has the same bit fields as an error frame containing an Active Error Flag. The key difference is in the cases that can trigger the transmission of an overload frame. Figure \ref{fig:twai-fields-of-an-overload-frame} below shows the bit fields of an overload frame.

\begin{figure}[H]
    \centering
    \fbox{\includegraphics[width=0.85\textwidth]{\modulefiles/figures/36-TWAI-fields-of-an-overload-frame__EN}}
    \caption{Fields of an Overload Frame}
    \label{fig:twai-fields-of-an-overload-frame}
\end{figure}

\input{\modulefiles/tables/36-TWAI-overload-frame__EN}

Overload frames will be transmitted in the following cases:

\begin{enumerate}
    \item A receiver requires a delay of the next data or remote frame.
    \item A dominant bit is detected at the first and second bit of intermission.
    \item A dominant bit is detected at the eighth (last) bit of an Error Delimiter. Note that in this case, TEC and REC will not be incremented (see Section \ref{sec:twai-errors} for more details).
\end{enumerate}

Transmitting an overload frame due to one of the above cases must also satisfy the following rules:
\begin{itemize}
        \item The start of an overload frame due to case 1 is only allowed to be started at the first bit time of an expected intermission.
        \item The start of an overload frame due to case 2 and 3 is only allowed to be started one bit after detecting the dominant bit.
        \item For case 1, a maximum of two overload frames may be generated in order to delay the transmission of the next data or remote frame.
\end{itemize}

\subsubsection{Interframe Space}

The Interframe Space acts as a separator between frames. Data frames and remote frames must be separated from preceding frames by an Interframe Space, regardless of the preceding frame's type (data frame, remote frame, error frame, or overload frame). However, error frames and overload frames do not need to be separated from preceding frames.

Figure \ref{fig:twai-fields-within-interframe-space} shows the fields within an Interframe Space:

\begin{figure}[ht]
    \centering
    \fbox{
\includegraphics[width=0.7\textwidth]{\modulefiles/figures/36-TWAI-fields-within-an-interframe-space__EN}}
    \caption{The Fields within an Interframe Space}
    \label{fig:twai-fields-within-interframe-space}
\end{figure}
\input{\modulefiles/tables/36-TWAI-interframe-space__EN}

\subsection{TWAI Errors}
\label{sec:twai-errors}
\subsubsection{Error Types}

Bus Errors in TWAI are categorized into the following types:

{\bfseries Bit Error}\\
A Bit Error occurs when a node transmits a bit value (i.e., dominant or recessive) but detects an opposite bit (e.g., a dominant bit is transmitted but a recessive is detected). However, if the transmitted bit is recessive and is located in the Arbitration Field or ACK Slot or Passive Error Flag, then detecting a dominant bit will not be considered as a Bit Error.

{\bfseries Stuff Error}\\
A stuff error occurs when six consecutive bits of the same value are detected (which violates the bit-stuffing encoding rules).

{\bfseries CRC Error}\\
A receiver of a data or remote frame will calculate CRC based on the bits it has received. A CRC error occurs when the CRC calculated by the receiver does not match the CRC sequence in the received data or remote Frame.

{\bfseries Format Error}\\
A Format Error occurs when a format-fixed bit field of a message contains an illegal bit. For example, the r1 and r0 fields must be dominant.

{\bfseries ACK Error}\\
An ACK Error occurs when a transmitter does not detect a dominant bit at the ACK Slot.

\subsubsection{Error States}

TWAI nodes implement fault confinement by maintaining two error counters in each node, where the counter values determine the error state. The two error counters are known as the Transmit Error Counter (TEC) and Receive Error Counter (REC). TWAI has the following error states:

{\bfseries Error Active}\\
An Error Active node is able to participate in bus communication and transmit an Active Error Flag when it detects an error.

{\bfseries Error Passive}\\
An Error Passive node is able to participate in bus communication and transmit an Passive Error Flag when it detects an error. Error Passive nodes that have transmitted a data or remote frame must also include the Suspend Transmission field in the subsequent Interframe Space.

{\bfseries Bus Off}\\
A Bus Off node is not permitted to influence the bus in any way (i.e., is not allowed to transmit data).

\subsubsection{Error Counters}

The TEC and REC are incremented/decremented according to the following rules. \textbf{Note that more than one rule can apply to a given message transfer.}

\begin{enumerate}
    \item When a receiver detects an error, the REC is increased by 1, except when the detected error was a Bit Error during the transmission of an Active Error Flag or an Overload Flag.
    \item When a receiver detects a dominant bit as the first bit after sending an Error Flag, the REC is increased by 8.
    \item When a transmitter sends an Error Flag, the TEC is increased by 8. However, the following scenarios are exempt from this rule:
    \begin{itemize}
        \item A transmitter is Error Passive and no dominant bit is detected when an Acknowledgment Error is detected and the Passive Error Flag is sent. In this case, the TEC should not be increased.
        \item A transmitter transmits an Error Flag due to a Stuff Error during Arbitration. If the stuffed bit should have been recessive but was monitored as dominant, then the TEC should not be increased.
    \end{itemize}
    \item If a transmitter detects a Bit Error whilst sending an Active Error Flag or Overload Flag, the TEC is increased by 8.
    \item If a receiver detects a Bit Error while sending an Active Error Flag or Overload Flag, the REC is increased by 8.
    \item A node can tolerate up to 7 consecutive dominant bits after sending an Active/Passive Error Flag, or Overload Flag. After detecting the 14th consecutive dominant bit when sending an Active Error Flag or Overload Flag, or the 8th consecutive dominant bit following a Passive Error Flag, a transmitter will increase its TEC by 8 and a receiver will increase its REC by 8. Every additional 8 consecutive dominant bits will also increase the TEC for transmitters or REC for receivers by 8 as well.
    \item When a transmitter has transmitted a message, which means getting ACK and no errors until the EOF is completed, the TEC is decremented by 1, unless the TEC is already at 0.
    \item When a receiver successfully receives a message, which means getting no errors before ACK Slot and successfully sending ACK, the REC is decremented accordingly.
    \begin{itemize}
        \item If the REC is between 1 and 127, the REC will be decremented by 1.
        \item If the REC is greater than 127, the REC will be set to 127.
        \item If the REC is 0, the REC will remain 0.
    \end{itemize}
    \item A node becomes Error Passive when its TEC and/or REC is greater than or equal to 128. Though the node becomes Error Passive, it still sends an Active Error Flag. Note that once the REC has reached to 128, any further increases to its value are invalid until the REC returns to a value less than 128.
    \item A node becomes Bus Off when its TEC is greater than or equal to 256.
    \item An Error Passive node becomes Error Active when both the TEC and REC are less than or equal to 127.
    \item A Bus Off node can become Error Active (with both its TEC and REC reset to 0) after it monitors 128 occurrences of 11 consecutive recessive bits on the bus.
\end{enumerate}

\subsection{TWAI Bit Timing}
\subsubsection{Nominal Bit}

The TWAI protocol allows a TWAI bus to operate at a particular bit rate. However, all nodes within a TWAI bus must operate at the same bit rate.

\begin{itemize}
    \item \textbf{The Nominal Bit Rate} is defined as the number of bits transmitted per second.
    \item \textbf{The Nominal Bit Time} is defined as \textbf{1/Nominal Bit Rate}.
\end{itemize}

A single Nominal Bit Time is divided into multiple segments, and each segment is made up of multiple Time Quanta. A \textbf{Time Quantum} is a minimum unit of time, and is implemented as some form of prescaled clock signal in each node. Figure \ref{fig:twai-bit-timing} illustrates the segments within a single Nominal Bit Time.

TWAI controllers will operate in time steps of one Time Quanta where the state of the TWAI bus is analyzed. If the bus states in two consecutive Time Quantas are different (i.e., recessive to dominant or vice versa), it means an edge is generated. The intersection of PBS1 and PBS2 is considered the Sample Point and the sampled bus value is considered the value of that bit.

\begin{figure}[ht]
    \centering
    \fbox{\includegraphics[width=0.55\textwidth]{\modulefiles/figures/36-TWAI-bit-timing__EN}}
    \caption{Layout of a Bit}
    \label{fig:twai-bit-timing}
\end{figure}

\input{\modulefiles/tables/36-TWAI-seg-of-a-nominal-bit-time__EN}

\subsubsection{Hard Synchronization and Resynchronization}

Due to clock skew and jitter, the bit timing of nodes on the same bus may become out of phase. Therefore, a bit edge may come before or after the SS. To ensure that the internal bit timing clocks of each node are kept in phase, TWAI has various methods of synchronization. The \textbf{Phase Error “e”} is measured in the number of Time Quanta and relative to the SS.

\begin{itemize}
    \item A positive Phase Error (e > 0) is when the edge lies after the SS and before the Sample Point (i.e., the edge is late).
    \item A negative Phase Error (e < 0) is when the edge lies after the Sample Point of the previous bit and before SS (i.e., the edge is early).
\end{itemize}

To correct for Phase Errors, there are two forms of synchronization, known as \textbf{Hard Synchronization} and \textbf{Resynchronization}. \textbf{Hard Synchronization} and \textbf{Resynchronization} obey the following rules:

\begin{itemize}
    \item Only one synchronization may occur in a single bit time.
    \item Synchronizations only occurs on recessive to dominant edges.
\end{itemize}

{\bfseries Hard Synchronization}\\
Hard Synchronization occurs on the recessive to dominant (i.e., the first SOF bit after Bus Idle) edges when the bus is idle. All nodes will restart their internal bit timings so that the recessive to dominant edge lies within the SS of the restarted bit timing.

{\bfseries Resynchronization}\\
Resynchronization occurs on recessive to dominant edges when the bus is not idle. If the edge has a positive Phase Error (e > 0), PBS1 is lengthened by a certain number of Time Quanta. If the edge has a negative Phase Error (e < 0), PBS2 will be shortened by a certain number of Time Quanta.

The number of Time Quanta to lengthen or shorten depends on the magnitude of the Phase Error, and is also limited by the Synchronization Jump Width (SJW) value which is programmable.

\begin{itemize}
    \item When the magnitude of the Phase Error (\textbf{e}) is less than or equal to the SJW, PBS1/PBS2 are lengthened/shortened by the \textbf{e} number of Time Quanta. This has a same effect as Hard Synchronization.
    \item When the magnitude of the Phase Error is greater to the SJW, PBS1/PBS2 are lengthened/shortened by the SJW number of Time Quanta. This means it may take multiple bits of synchronization before the Phase Error is entirely corrected.
\end{itemize}

\section{Architectural Overview}

\begin{figure}[ht]
    \centering
    \fbox{\includegraphics[width=0.8\textwidth]{\modulefiles/figures/36-TWAI-overview.eps}}
    \caption{TWAI Overview Diagram}
    \label{fig:twai-overview-diagram-diagram}
\end{figure}

The major functional blocks of the TWAI controller are shown in Figure \ref{fig:twai-overview-diagram-diagram}.

\subsection{Registers Block}
The \chipname{} CPU accesses peripherals using 32-bit aligned words. However, the majority of registers in the TWAI controller only contain useful data at the least significant byte (bits [7:0]). Therefore, in these registers, bits [31:8] are ignored on writes, and return 0 on reads.

{\bfseries Configuration Registers}\\
The configuration registers store various configuration items for the TWAI controller such as bit rates, operation mode, Acceptance Filter, etc. Configuration registers can only be modified whilst the TWAI controller is in Reset Mode (See Section \ref{sec:twai-modes}).

{\bfseries Command Registers}\\
The command register is used by the CPU to drive the TWAI controller to initiate certain actions such as transmitting a message or clearing the Receive Buffer. The command register can only be modified when the TWAI controller is in Operation Mode (see section \ref{sec:twai-modes}).

{\bfseries Interrupt \verb+&+ Status Registers}\\
The interrupt register indicates what events have occurred in the TWAI controller (each event is represented by a separate bit). The status register indicates the current status of the TWAI controller.

{\bfseries Error Management Registers}\\
The error management registers include error counters and capture registers. The error counter registers represent TEC and REC values. The capture registers will record information about instances where TWAI controller detects a bus error, or when it loses arbitration.

{\bfseries Transmit Buffer Registers}\\
The transmit buffer is a 13-byte buffer used to store a TWAI message to be transmitted.

{\bfseries Receive Buffer Registers}\\
The Receive Buffer is a 13-byte buffer which stores a single message. The Receive Buffer acts as a window of Receive FIFO, whose first message will be mapped into the Receive Buffer.

Note that the Transmit Buffer registers, Receive Buffer registers, and the Acceptance Filter registers share the same address range (offset 0x0040 to 0x0070). Their access is governed by the following rules:

\begin{itemize}
    \item When the TWAI controller is in Reset Mode, all reads and writes to the address range maps to the Acceptance Filter registers.
    \item When the TWAI controller is in Operation Mode:
    \begin{itemize}
        \item All reads to the address range maps to the Receive Buffer registers.
        \item All writes to the address range maps to the Transmit Buffer registers.
    \end{itemize}
\end{itemize}

\subsection{Bit Stream Processor}
The Bit Stream Processing (BSP) module handles the frame processing of data from the Transmit Buffer (e.g. bit stuffing and additional CRC fields) and generates a bit stream for the Bit Timing Logic (BTL) module. At the same time, the BSP module is also responsible for processing the received bit stream (e.g., de-stuffing and verifying CRC) from the BTL module and placing the message into the Receive FIFO. The BSP will also detect errors on the TWAI bus and report them to the Error Management Logic (EML).

\subsection{Error Management Logic}
The Error Management Logic (EML) module updates the TEC and REC, records error information like error types and positions, and updates the error state of the TWAI controller such that the BSP module generates the correct Error Flags. Furthermore, this module also records the bit position when the TWAI controller loses arbitration.

\subsection{Bit Timing Logic}
The Bit Timing Logic (BTL) module transmits and receives messages at the configured bit rate. The BTL module also handles bit timing synchronization so that communication remains stable. A single bit time consists of multiple programmable segments that allows users to set the length of each segment to account for factors such as propagation delay and controller processing time, etc.

\subsection{Acceptance Filter}
The Acceptance Filter is a programmable message filtering unit that allows the TWAI controller to accept or reject a received message based on the message’s ID field. Only accepted messages will be stored in the Receive FIFO. The Acceptance Filter’s registers can be programmed to specify a single filter, or two separate filters (dual filter mode).

\subsection{Receive FIFO}
The Receive FIFO is a 64-byte buffer (inside the TWAI controller) that stores received messages accepted by the Acceptance Filter. Messages in the Receive FIFO can vary in size (between 3 to 13 bytes). When the Receive FIFO is full (or does not have enough space to store the current received message in its entirety), the Overrun Interrupt will be triggered, and any subsequent received messages will be lost until adequate space is cleared in the Receive FIFO. The first message in the Receive FIFO will be mapped to the 13-byte Receive Buffer until that message is cleared (using the Release Receive Buffer command bit). After being cleared, the Receive Buffer will map to the next message in the Receive FIFO, and the space occupied by the previous message in the Receive FIFO can be used to receive new messages.

\section{Functional Description}

\subsection{Modes}
\label{sec:twai-modes}
The \chipname{} TWAI controller has two working modes: Reset Mode and Operation Mode. Reset Mode and Operation Mode are entered by setting or clearing the \hyperref[fielddesc:TWAIRESETMODE]{TWAI\_RESET\_MODE} bit.

\subsubsection{Reset Mode}
Entering Reset Mode is required in order to modify the various configuration registers of the TWAI controller. When entering Reset Mode, the TWAI controller is essentially disconnected from the TWAI bus. When in Reset Mode, the TWAI controller will not be able to transmit any messages (including error signals). Any transmission in progress is immediately terminated. Likewise, the TWAI controller will not be able to receive any messages either.

\subsubsection{Operation Mode}
In operation mode, the TWAI controller connects to the bus and write-protect all configuration registers to ensure consistency during operation. When in Operation Mode, the TWAI controller can transmit and receive messages (including error signaling) depending on which operation sub-mode the TWAI controller was configured with. The TWAI controller supports the following operation sub-modes:

\begin{itemize}
    \item {\bfseries Normal Mode:} The TWAI controller can transmit and receive messages including error signals (such as error and overload Frames).
    \item {\bfseries Self-test Mode:} Self-test mode is similar to normal Mode, but the TWAI controller will consider the transmission of a data or remote frame successful and do not generate an ACK error even if it was not acknowledged. This is commonly used when the TWAI controller does self-test.
    \item {\bfseries Listen-only Mode:} The TWAI controller will be able to receive messages, but will remain completely passive on the TWAI bus. Thus, the TWAI controller will not be able to transmit any messages, acknowledgments, or error signals. The error counters will remain frozen. This mode is useful for TWAI bus monitoring.
\end{itemize}

Note that when exiting Reset Mode (i.e., entering Operation Mode), the TWAI controller must wait for 11 consecutive recessive bits to occur before being able to fully connect the TWAI bus (i.e., be able to transmit or receive).

\subsection{Bit Timing}
The operating bit rate of the TWAI controller must be configured whilst the TWAI controller is in Reset Mode. The bit rate is configured using \hyperref[regdesc:TWAIBUSTIMING0REG]{TWAI\_BUS\_TIMING\_0\_REG} and \hyperref[regdesc:TWAIBUSTIMING1REG]{TWAI\_BUS\_TIMING\_1\_REG}, and the two registers contain the following fields:

The following Table \ref{tab:twai-bit-info-bus-timing-0-reg} illustrates the bit fields of \hyperref[regdesc:TWAIBUSTIMING0REG]{TWAI\_BUS\_TIMING\_0\_REG}. The frequency of the TWAI core clock has multiple clock sources that can be configured by the user as needed. See Chapter \ref{mod:resclk} \textit{\nameref{mod:resclk}} for detailed configuration instructions.

\begin{longtable}[c]{|p{2cm}|p{1.5cm}|p{1.5cm}|p{1.5cm}|p{1.5cm}|p{1.5cm}|p{1.5cm}|p{1.5cm}|}
\caption{Bit Information of \hyperref[regdesc:TWAIBUSTIMING0REG]{TWAI\_BUS\_TIMING\_0\_REG} (0x18)}
\label{tab:twai-bit-info-bus-timing-0-reg}\\\hline
\rowcolor{lightgray}
  \textbf{Bit 31-16} & \textbf{Bit 15} & \textbf{Bit 14} & \textbf{Bit 13} & \textbf{Bit 12} & \textbf{......} & \textbf{Bit 1} & \textbf{Bit 0}\\\hline 
  Reserved& SJW.1 & SJW.0 & BRP.13 & BRP.12 & ...... & BRP.1 & BRP.0\\\hline
\end{longtable}
\vspace{-1.5em}
{\bfseries Notes:}\\
\begin{itemize}
    \item BRP: The TWAI Time Quanta clock is derived from the XTAL clock (the default is 40 MHz and is configured). The Baud Rate Prescaler (BRP) field is used to define the prescaler according to the equation below, where t$_{Tq}$ is the Time Quanta clock cycle and t$_{CLK}$ is TWAI core clock cycle:\\

t$_{Tq}$ = 2 × t$_{CLK}$ × (2$^{13}$ × BRP.13 + 2$^{12}$ × BRP.12 + 2$^{11}$ × BRP.11 + ... + 2$^1$ × BRP.1 + 2$^0$ × BRP.0 + 1)
    \item SJW: Synchronization Jump Width (SJW) is configured in SJW.0 and SJW.1 where SJW = (2 x SJW.1 + SJW.0 + 1).
\end{itemize}

The following Table \ref{tab:twai-bit-info-bus-timing-1-reg} illustrates the bit fields of \hyperref[regdesc:TWAIBUSTIMING1REG]{TWAI\_BUS\_TIMING\_1\_REG}.

\begin{longtable}[c]{|p{1.4cm}|p{1.4cm}|p{1.4cm}|p{1.4cm}|p{1.4cm}|p{1.4cm}|p{1.4cm}|p{1.4cm}|p{1.4cm}|}
    \caption{Bit Information of \hyperref[regdesc:TWAIBUSTIMING1REG]{TWAI\_BUS\_TIMING\_1\_REG} (0x1c)}
    \label{tab:twai-bit-info-bus-timing-1-reg}\\
    \hline
       \rowcolor{lightgray}
   \textbf{Bit 31-8} & \textbf{Bit 7} & \textbf{Bit 6} & \textbf{Bit 5} & \textbf{Bit 4} & \textbf{Bit 3} & \textbf{Bit 2} & \textbf{Bit 1} & \textbf{Bit 0}\\\hline 
   Reserved & SAM & PBS2.2 & PBS2.1 & PBS2.0 & PBS1.3 &PBS1.2 & PBS1.1 & PBS1.0\\
    \hline
\end{longtable}
\vspace{-1.5em}
{\bfseries Notes:}\\
\begin{itemize}
    \item PBS1: The number of Time Quanta in Phase Buffer Segment 1 is defined according to the following equation: (8 x PBS1.3 + 4 x PBS1.2 + 2 x PBS1.1 + PBS1.0 + 1)。
    \item PBS2: The number of Time Quanta in Phase Buffer Segment 2 is defined according to the following equation: (4 x PBS2.2 + 2 x PBS2.1 + PBS2.0 + 1)。
    \item SAM: Enables triple sampling if set to 1. This is useful for low/medium speed buses to filter spikes on the bus line.
\end{itemize}

\subsection{Interrupt Management}
\label{sec:twai-interrupt-management}
The \chipname{} TWAI controller provides eight interrupts, each represented by a single bit in the \hyperref[regdesc:TWAIINTSTREG]{TWAI\_INT\_ST\_REG}. For a particular interrupt to be triggered, the corresponding enable bit in \hyperref[regdesc:TWAIINTENAREG]{TWAI\_INT\_ENA\_REG} must be set.

The TWAI controller provides the following interrupts:

\begin{itemize}
    \item Receive Interrupt
    \item Transmit Interrupt
    \item Error Warning Interrupt
    \item Data Overrun Interrupt
    \item Error Passive Interrupt
    \item Arbitration Lost Interrupt
    \item Bus Error Interrupt
    \item Bus Idle Status Interrupt
\end{itemize}

The TWAI controller’s interrupt signal to the interrupt matrix will be asserted whenever one or more interrupt bits are set in the \hyperref[regdesc:TWAIINTSTREG]{TWAI\_INT\_ST\_REG}, and de-asserted when all bits in \hyperref[regdesc:TWAIINTSTREG]{TWAI\_INT\_ST\_REG} are cleared. The majority of interrupt bits in \hyperref[regdesc:TWAIINTSTREG]{TWAI\_INT\_ST\_REG} are automatically cleared when the register is read, except for the Receive Interrupt which can only be cleared when all the messages are released by setting the \hyperref[fielddesc:TWAIRELEASEBUF]{TWAI\_RELEASE\_BUF} bit.

\subsubsection{Receive Interrupt (RXI)}
The Receive Interrupt (RXI) is asserted whenever the TWAI controller has received messages that are pending to be read from the Receive Buffer (i.e., when \hyperref[regdesc:TWAIRXMESSAGECNTREG]{TWAI\_RX\_MESSAGE\_CNT\_REG} > 0). Pending received messages includes valid messages in the Receive FIFO and also overrun messages. The RXI will not be deasserted until all pending received messages are cleared using the \hyperref[fielddesc:TWAIRELEASEBUF]{TWAI\_RELEASE\_BUF} command bit.


\subsubsection{Transmit Interrupt (TXI)}
The Transmit Interrupt (TXI) is triggered whenever Transmit Buffer becomes free, indicating another message can be loaded into the Transmit Buffer to be transmitted. The Transmit Buffer becomes free under the following scenarios:

\begin{itemize}
    \item A message transmission has completed successfully, i.e., acknowledged without any errors. Any failed messages will automatically be resent.
    \item A single shot transmission has completed (successfully or unsuccessfully, indicated by the \hyperref[fielddesc:TWAITXCOMPLETE]{TWAI\_TX\_COMPLETE} bit).
    \item A message transmission was aborted using the \hyperref[fielddesc:TWAIABORTTX]{TWAI\_ABORT\_TX} command bit.
\end{itemize}

\subsubsection{Error Warning Interrupt (EWI)}
The Error Warning Interrupt (EWI) is triggered whenever there is a change to the \hyperref[fielddesc:TWAIERRST]{TWAI\_ERR\_ST} and \hyperref[fielddesc:TWAIBUSOFFST]{TWAI\_BUS\_OFF\_ST} bits of the \hyperref[regdesc:TWAISTATUSREG]{TWAI\_STATUS\_REG} (i.e., transition from 0 to 1 or vice versa). Thus, an EWI could indicate one of the following events, depending on the values \hyperref[fielddesc:TWAIERRST]{TWAI\_ERR\_ST} and \hyperref[fielddesc:TWAIBUSOFFST]{TWAI\_BUS\_OFF\_ST} at the moment when the EWI is triggered.

\begin{itemize}
    \item If \hyperref[fielddesc:TWAIERRST]{TWAI\_ERR\_ST} = 0 and \hyperref[fielddesc:TWAIBUSOFFST]{TWAI\_BUS\_OFF\_ST} = 0:
    \begin{itemize}
        \item If the TWAI controller was in the Error Active state, it indicates both the TEC and REC have returned below the threshold value set by \hyperref[regdesc:TWAIERRWARNINGLIMITREG]{TWAI\_ERR\_WARNING\_LIMIT\_REG}.
        \item If the TWAI controller was previously in the Bus Off Recovery state, it indicates that Bus Recovery has completed successfully.
    \end{itemize}
    \item If \hyperref[fielddesc:TWAIERRST]{TWAI\_ERR\_ST} = 1 and \hyperref[fielddesc:TWAIBUSOFFST]{TWAI\_BUS\_OFF\_ST} = 0: The TEC or REC error counters have exceeded the threshold value set by  \hyperref[regdesc:TWAIERRWARNINGLIMITREG]{TWAI\_ERR\_WARNING\_LIMIT\_REG}.
    \item If  \hyperref[fielddesc:TWAIERRST]{TWAI\_ERR\_ST} = 1 and \hyperref[fielddesc:TWAIBUSOFFST]{TWAI\_BUS\_OFF\_ST} = 1: The TWAI controller has entered the BUS\_OFF state (due to the TEC >= 256).
    \item If  \hyperref[fielddesc:TWAIERRST]{TWAI\_ERR\_ST} = 0 and \hyperref[fielddesc:TWAIBUSOFFST]{TWAI\_BUS\_OFF\_ST} = 1: The TWAI controller’s TEC has dropped below the threshold value set by \hyperref[regdesc:TWAIERRWARNINGLIMITREG]{TWAI\_ERR\_WARNING\_LIMIT\_REG} during BUS\_OFF recovery.
\end{itemize}

\subsubsection{Data Overrun Interrupt (DOI)}
The Data Overrun Interrupt (DOI) is triggered whenever the Receive FIFO has overrun. The DOI indicates that the Receive FIFO is full and should be cleared immediately to prevent any further overrun messages.

The DOI is only triggered by the first message that causes the Receive FIFO to overrun (i.e., the transition from the Receive FIFO not being full to the Receive FIFO overrunning). Any subsequent overrun messages will not trigger the DOI again. The DOI could be triggered again when all received messages (valid or overrun) have been cleared.


\subsubsection{Error Passive Interrupt (TXI)}
The Error Passive Interrupt (EPI) is triggered whenever the TWAI controller switches from Error Active to Error Passive, or vice versa.

\subsubsection{Arbitration Lost Interrupt (ALI)}
The Arbitration Lost Interrupt (ALI) is triggered whenever the TWAI controller is attempting to transmit a message and loses arbitration. The bit position where the TWAI controller lost arbitration is automatically recorded in Arbitration Lost Capture register (\hyperref[regdesc:TWAIARBLOSTCAPREG]{TWAI\_ARB\_LOST\_CAP\_REG}). When the ALI occurs again, the Arbitration Lost Capture register will no longer record new bit location until it is cleared (via CPU reading this register).

\subsubsection{Bus Error Interrupt (BEI)}
The Bus Error Interrupt (BEI) is triggered whenever TWAI controller detects an error on the TWAI bus. When a bus error occurs, the Bus Error type and its bit position are automatically recorded in the Error Code Capture register (\hyperref[regdesc:TWAIERRCODECAPREG]{TWAI\_ERR\_CODE\_CAP\_REG}). When the BEI occurs again, the Error Code Capture register will no longer record new error information until it is cleared (via a read from the CPU).

\subsubsection{Bus Idle Status Interrupt (BISI)}
The Bus Idle Status Interrupt (BISI) is triggered when the number of clock cycles of the TWAI controller in the idle status exceeds the pre-configured value in the TWAI\_IDLE\_INTR\_CNT\_REG register. Users can configure this interrupt to get the TWAI controller idle status and further decide whether to turn off the external TWAI receiver to reduce the overall power consumption (see Section \ref{sec:twai-transceiver-auto-standby}).

\subsection{Transmit and Receive Buffers}
\subsubsection{Overview of Buffers}

\input{\modulefiles/tables/36-TWAI-buffer-layout__EN}

Table \ref{tab:twai-buffer-layout-std-extended-frame-format} illustrates the layout of the Transmit Buffer and Receive Buffer registers. Both the Transmit and Receive Buffer registers share the same address space and are only accessible when the TWAI controller is in Operation Mode. The CPU accesses Transmit Buffer registers for write operations, and Receive Buffer registers for read operations. Both buffers share the exact same register layout and fields to store a message (received or to be transmitted).
The Transmit Buffer registers are used to configure a TWAI message to be transmitted. The CPU would write to the Transmit Buffer registers specifying the message’s frame type, frame format, frame ID, and frame data (payload). Once the Transmit Buffer is configured, the CPU would then initiate the transmission by setting the  \hyperref[fielddesc:TWAITXREQ]{TWAI\_TX\_REQ} bit in \hyperref[regdesc:TWAICMDREG]{TWAI\_CMD\_REG}.

\begin{itemize}
    \item For a self-reception request, set the \hyperref[fielddesc:TWAISELFRXREQ]{TWAI\_SELF\_RX\_REQ} bit instead.
    \item For a single-shot transmission, set both the \hyperref[fielddesc:TWAITXREQ]{TWAI\_TX\_REQ} and the \hyperref[fielddesc:TWAIABORTTX]{TWAI\_ABORT\_TX} simultaneously.
\end{itemize}

The Receive Buffer registers map the first message in the Receive FIFO. The CPU would read the Receive Buffer registers to obtain the first message’s frame type, frame format, frame ID, and frame data (payload). Once the message has been read from the Receive Buffer registers, the CPU can set the \hyperref[fielddesc:TWAIRELEASEBUF]{TWAI\_RELEASE\_BUF} bit in \hyperref[regdesc:TWAICMDREG]{TWAI\_CMD\_REG} to clear the Receive Buffer registers. If there are still messages in the Receive FIFO, the Receive Buffer registers will map the first of the remaining messages again.

\subsubsection{Frame Information}
The frame information is one byte long and specifies a message’s frame type, frame format, and length of data. The frame information fields are shown in Table \ref{tab:twai-tx-rx-frame-info-sff-eff}.

\begin{longtable}[c]{|p{1.4cm}|p{1.4cm}|p{1.4cm}|p{1.4cm}|p{1.4cm}|p{1.4cm}|p{1.4cm}|p{1.4cm}|p{1.4cm}|}
    \caption{TX/RX Frame Information (SFF/EFF)；TWAI Address 0x40}
    \label{tab:twai-tx-rx-frame-info-sff-eff}\\
    \hline
       \rowcolor{lightgray}
    \textbf{Bit 31-8} & \textbf{Bit 7} & \textbf{Bit 6} & \textbf{Bit 5} & \textbf{Bit 4} & \textbf{Bit 3} & \textbf{Bit 2} & \textbf{Bit 1} & \textbf{Bit 0}\\\hline
    Reserved & FF$^1$ & RTR$^2$  & X$^3$  & X$^3$  & DLC.3$^4$ & DLC.2$^4$ & DLC.1$^4$ & DLC.0$^4$\\
    \hline
\end{longtable}
\vspace{-1.5em}
{\bfseries Notes:}\\
\begin{enumerate}
    \item FF: The Frame Format (FF) bit specifies whether the message is Extended Frame Format (EFF) or Standard Frame Format (SFF). The message is EFF when FF bit is 1, and SFF when FF bit is 0.
    \item RTR: The Remote Transmission Request (RTR) bit specifies whether the message is a data frame or a remote frame. The message is a remote frame when the RTR bit is 1, and a data frame when the RTR bit is 0.
    \item X: Don’t care, can be any value.
    \item DLC: The Data Length Code (DLC) field specifies the number of data bytes for a data frame, or the number of data bytes to request in a remote frame. TWAI data frames are limited to a maximum payload of 8 data bytes, and thus the DLC should range anywhere from 0 to 8.
\end{enumerate}


\subsubsection{Frame Identifier}
The Frame Identifier fields occupies two-byte (11-bit) long if the message is SFF, and four-byte (29-bit) long if the message is EFF.

The Frame Identifier fields for an SFF (11-bit) message is shown in Table \ref{tab:twai-tx-rx-identifier-1-sff} \verb+~+ \ref{tab:twai-tx-rx-identifier-2-sff}. 

\begin{longtable}[c]{|p{1.4cm}|p{1.4cm}|p{1.4cm}|p{1.4cm}|p{1.4cm}|p{1.4cm}|p{1.4cm}|p{1.4cm}|p{1.4cm}|}
    \caption{TX/RX Identifier 1 (SFF); TWAI Address 0x44}
    \label{tab:twai-tx-rx-identifier-1-sff}\\
    \hline
    \rowcolor{lightgray}
    \textbf{Bit 31-8} & \textbf{Bit 7} & \textbf{Bit 6} & \textbf{Bit 5} & \textbf{Bit 4} & \textbf{Bit 3} & \textbf{Bit 2} & \textbf{Bit 1} & \textbf{Bit 0}\\\hline
    
    Reserved & ID.10 & ID.9 & ID.8 & ID.7 & ID.6 & ID.5 & ID.4 & ID.3\\
    \hline
\end{longtable}

\begin{longtable}[c]{|p{1.4cm}|p{1.4cm}|p{1.4cm}|p{1.4cm}|p{1.4cm}|p{1.4cm}|p{1.4cm}|p{1.4cm}|p{1.4cm}|}
    \caption{TX/RX Identifier 2 (SFF); TWAI Address 0x48}
    \label{tab:twai-tx-rx-identifier-2-sff}\\
    \hline
       \rowcolor{lightgray}
    \textbf{Bit 31-8} & \textbf{Bit 7} & \textbf{Bit 6} & \textbf{Bit 5} & \textbf{Bit 4} & \textbf{Bit 3} & \textbf{Bit 2} & \textbf{Bit 1} & \textbf{Bit 0}\\\hline
    Reserved & ID.2 & ID.1 & ID.0 & X$^1$ & X$^2$ & X$^2$ & X$^2$ & X$^2$\\
    \hline
\end{longtable}
\vspace{-1.5em}
{\bfseries Notes:}\\
\begin{enumerate}
    \item Don’t care. Recommended to be compatible with receive buffer (i.e., set to RTR ) in case of using the self reception functionality (or together with self-test functionality).
    \item Don’t care. Recommended to be compatible with receive buffer (i.e., set to 0 ) in case of using the self reception functionality (or together with self-test functionality). 
\end{enumerate}


The Frame Identifier fields for an EFF (29-bits) message is shown in Table \ref{tab:twai-tx-rx-identifier-1-eff} \verb+~+ \ref{tab:twai-tx-rx-identifier-4-eff}.

\begin{longtable}[c]{|p{1.4cm}|p{1.4cm}|p{1.4cm}|p{1.4cm}|p{1.4cm}|p{1.4cm}|p{1.4cm}|p{1.4cm}|p{1.4cm}|}
    \caption{TX/RX Identifier 1 (EFF); TWAI Address 0x44}
    \label{tab:twai-tx-rx-identifier-1-eff}\\
    \hline
       \rowcolor{lightgray}
    \textbf{Bit 31-8} & \textbf{Bit 7} & \textbf{Bit 6} & \textbf{Bit 5} & \textbf{Bit 4} & \textbf{Bit 3} & \textbf{Bit 2} & \textbf{Bit 1} & \textbf{Bit 0}\\\hline
    Reserved & ID.28 & ID.27 & ID.26 & ID.25 & ID.24 & ID.23 & ID.22 & ID.21\\
    \hline
\end{longtable}

\begin{longtable}[c]{|p{1.4cm}|p{1.4cm}|p{1.4cm}|p{1.4cm}|p{1.4cm}|p{1.4cm}|p{1.4cm}|p{1.4cm}|p{1.4cm}|}
    \caption{TX/RX Identifier 2 (EFF); TWAI Address 0x48}
    \label{tab:twai-tx-rx-identifier-2-eff}\\
    \hline
       \rowcolor{lightgray}
    \textbf{Bit 31-8} & \textbf{Bit 7} & \textbf{Bit 6} & \textbf{Bit 5} & \textbf{Bit 4} & \textbf{Bit 3} & \textbf{Bit 2} & \textbf{Bit 1} & \textbf{Bit 0}\\\hline
    Reserved & ID.20 & ID.19 & ID.18 & ID.17 & ID.16 & ID.15 & ID.14 & ID.13\\
    \hline
\end{longtable}

\begin{longtable}[c]{|p{1.4cm}|p{1.4cm}|p{1.4cm}|p{1.4cm}|p{1.4cm}|p{1.4cm}|p{1.4cm}|p{1.4cm}|p{1.4cm}|}
    \caption{TX/RX Identifier 3 (EFF); TWAI Address 0x4c}
    \label{tab:twai-tx-rx-identifier-3-eff}\\
    \hline
       \rowcolor{lightgray}
    \textbf{Bit 31-8} & \textbf{Bit 7} & \textbf{Bit 6} & \textbf{Bit 5} & \textbf{Bit 4} & \textbf{Bit 3} & \textbf{Bit 2} & \textbf{Bit 1} & \textbf{Bit 0}\\\hline 
    Reserved & ID.12 & ID.11 & ID.10 & ID.9 & ID.8 & ID.7 & ID.6 & ID.5\\
    \hline
\end{longtable}

\begin{longtable}[c]{|p{1.4cm}|p{1.4cm}|p{1.4cm}|p{1.4cm}|p{1.4cm}|p{1.4cm}|p{1.4cm}|p{1.4cm}|p{1.4cm}|}
    \caption{TX/RX Identifier 4 (EFF); TWAI Address 0x50}
    \label{tab:twai-tx-rx-identifier-4-eff}\\
    \hline
       \rowcolor{lightgray}
    \textbf{Bit 31-8} & \textbf{Bit 7} & \textbf{Bit 6} & \textbf{Bit 5} & \textbf{Bit 4} & \textbf{Bit 3} & \textbf{Bit 2} & \textbf{Bit 1} & \textbf{Bit 0}\\\hline
    Reserved & ID.4 & ID.3 & ID.2 & ID.1 & ID.0 & X$^1$ & X$^2$ & X$^2$\\
    \hline
\end{longtable}
\vspace{-1.5em}
{\bfseries Notes:}\\
\begin{enumerate}
    \item Don’t care. Recommended to be compatible with receive buffer (i.e., set to RTR ) in case of using the self reception functionality (or together with self-test functionality).
    \item Don’t care. Recommended to be compatible with receive buffer (i.e., set to 0 ) in case of using the self reception functionality (or together with self-test functionality). 
\end{enumerate}

\subsubsection{Frame Data}
The Frame Data field contains the payloads of transmitted or received data frame, and can range from 0 to 8 bytes. The number of valid bytes should be equal to the DLC. However, if the DLC is larger than eight bytes, the number of valid bytes would still be limited to eight. Remote frames do not have data payloads, so their Frame Data fields will be unused.

For example, when transmitting a data frame with five bytes, the CPU should write five to the DLC field, and then write data to the corresponding register of the first to the fifth data field. Likewise, when the CPU receives a data frame with a DLC of five data bytes, only the first to the fifth data byte will contain valid payload data for the CPU to read.

\subsection{Receive FIFO and Data Overruns}

The Receive FIFO is a 64-byte internal buffer used to store received messages in First In First Out order. A single received message can occupy between 3 to 13 bytes of space in the Receive FIFO, and their endianness is identical to the register layout of the Receive Buffer registers. The Receive Buffer registers are mapped to the bytes of the first message in the Receive FIFO.

When the TWAI controller receives a message, it will increment the value of \hyperref[fielddesc:TWAIRXMESSAGECOUNTER]{TWAI\_RX\_MESSAGE\_COUNTER} by 1, with a maximum of 64. If there is adequate space in the Receive FIFO, the message contents will be written into the Receive FIFO. Once a message has been read from the Receive Buffer, the \hyperref[fielddesc:TWAIRELEASEBUF]{TWAI\_RELEASE\_BUF} bit should be set. This will decrement \hyperref[fielddesc:TWAIRXMESSAGECOUNTER]{TWAI\_RX\_MESSAGE\_COUNTER} by 1 and free the space occupied by the first message in the Receive FIFO. The Receive Buffer will then map to the next message in the Receive FIFO.

A data overrun occurs when the TWAI controller receives a message, but the Receive FIFO lacks adequate free space to store the received message in its entirety (either due to the message contents being larger than the free space in the Receive FIFO, or the Receive FIFO being completely full). 

When a data overrun occurs:
\begin{itemize}
    \item The free space left in the Receive FIFO is filled with the partial contents of the overrun message. If the Receive FIFO is already full, then none of the overrun message’s contents will be stored.
    \item When data in the Receive FIFO overruns for the first time, a Data Overrun Interrupt will be triggered.
    \item Each overrun message will still increment the  \hyperref[fielddesc:TWAIRXMESSAGECOUNTER]{TWAI\_RX\_MESSAGE\_COUNTER} up to a maximum of 64.
    \item The Receive FIFO will internally mark overrun messages as invalid. The \hyperref[fielddesc:TWAIMISSST]{TWAI\_MISS\_ST} bit can be used to determine whether the message currently mapped to by the Receive Buffer is valid or overrun.
\end{itemize}

To clear an overrun Receive FIFO, the \hyperref[fielddesc:TWAIRELEASEBUF]{TWAI\_RELEASE\_BUF} must be called repeatedly until \\\hyperref[fielddesc:TWAIRXMESSAGECOUNTER]{TWAI\_RX\_MESSAGE\_COUNTER} is 0. This requires users to read all valid messages in the Receive FIFO and clear all overrun messages.

\subsection{Acceptance Filter}

The Acceptance Filter allows the TWAI controller to filter out received messages based on their ID (and optionally their first data byte and frame type). Only accepted messages are passed on to the Receive FIFO. The use of Acceptance Filters allows a more lightweight operation of the TWAI controller (e.g., less use of Receive FIFO, fewer Receive Interrupts) since the TWAI Controller only need to handle a subset of messages.

The Acceptance Filter configuration registers can only be accessed whilst the TWAI controller is in Reset Mode, since they share the same address spaces with the Transmit Buffer and Receive Buffer registers.

The configuration registers consist of a 32-bit Acceptance Code Value and a 32-bit Acceptance Mask Value. The Acceptance Code value specifies a bit pattern which each filtered bit of the message must match in order for the message to be accepted. The Acceptance Mask Value is able to mask out certain bits of the Code value (i.e., set as “Don't Care” bits). Each filtered bit of the message must either match the acceptance code or be masked in order for the message to be accepted, as demonstrated in Figure \ref{fig:twai-acceptance-filter}.

\begin{figure}[ht]
    \centering
    \fbox{\includegraphics[width=1.0\textwidth]{\modulefiles/figures/36-TWAI-acceptance-filter}}
    \caption{Acceptance Filter}
    \label{fig:twai-acceptance-filter}
\end{figure}

The TWAI controller Acceptance Filter allows the 32-bit Acceptance Code and Mask Values to either define a single filter (i.e., Single Filter Mode), or two filters (i.e., Dual Filter Mode). How the Acceptance Filter interprets the 32-bit code and mask values is dependent on filter mode and the format of received messages (i.e., SFF or EFF).

\subsubsection{Single Filter Mode}

Single Filter Mode is enabled by setting the \hyperref[fielddesc:TWAIRXFILTERMODE]{TWAI\_RX\_FILTER\_MODE} bit to 1. This will cause the 32-bit code and mask values to define a single filter. The single filter can filter the following bits of a data or remote frame:

\begin{itemize}
    \item SFF
    \begin{itemize}
        \item The entire 11-bit ID
        \item RTR bit
        \item Data byte 1 and Data byte 2
    \end{itemize}
    \item EFF
    \begin{itemize}
        \item The entire 29-bit ID
        \item RTR bit
    \end{itemize}
\end{itemize}

The following Figure \ref{fig:twai-single-filter-modee} illustrates how the 32-bit code and mask values will be interpreted under Single Filter Mode.

\begin{figure}[ht]
    \centering
    \fbox{\includegraphics[width=0.8\textwidth]{\modulefiles/figures/36-TWAI-single-filter-mode__EN}}
    \caption{Single Filter Mode}
    \label{fig:twai-single-filter-modee}
\end{figure}

\subsubsection{Dual Filter Mode}

Dual Filter Mode is enabled by clearing the \hyperref[fielddesc:TWAIRXFILTERMODE]{TWAI\_RX\_FILTER\_MODE} bit to 0. This will cause the 32-bit code and mask values to define a two separate filters referred to as filter 1 or filter 2. Under Dual Filter Mode, a message will be accepted if it is accepted by one of the two filters.

The two filters can filter the following bits of a data or remote frame:

\begin{itemize}
    \item SFF
    \begin{itemize}
        \item The entire 11-bit ID
        \item RTR bit
        \item Data byte 1 (for filter 1 only)
    \end{itemize}
    \item EFF
    \begin{itemize}
        \item The first 16 bits of the 29-bit ID
    \end{itemize}
\end{itemize}

The following Figure \ref{fig:twai-dual-filter-modee} illustrates how the 32-bit code and mask values will be interpreted in Dual Filter Mode.

\begin{figure}[ht!]
    \centering
    \fbox{\includegraphics[width=0.6\textwidth]{\modulefiles/figures/36-TWAI-dual-filter-mode__CN}}
    \caption{Dual Filter Mode}
    \label{fig:twai-dual-filter-modee}
\end{figure}

\subsection{Error Management}

The TWAI protocol requires that each TWAI node maintains the Transmit Error Counter (TEC) and Receive Error Counter (REC). The value of both error counters determines the current error state of the TWAI controller (i.e., Error Active, Error Passive, Bus-Off). The TWAI controller stores the TEC and REC values in \hyperref[regdesc:TWAITXERRCNTREG]{TWAI\_TX\_ERR\_CNT\_REG} and \hyperref[regdesc:TWAIRXERRCNTREG]{TWAI\_RX\_ERR\_CNT\_REG} respectively, and they can be read by the CPU anytime. In addition to the error states, the TWAI controller also offers an Error Warning Limit (EWL) feature that can warn users of the occurrence of severe bus errors before the TWAI controller enters the Error Passive state.

The current error state of the TWAI controller is indicated via a combination of the following values and status bits: TEC, REC, \hyperref[fielddesc:TWAIERRST]{TWAI\_ERR\_ST}, and \hyperref[fielddesc:TWAIBUSOFFST]{TWAI\_BUS\_OFF\_ST}. Certain changes to these values and bits will also trigger interrupts, thus allowing the users to be notified of error state transitions (see section \ref{sec:twai-interrupt-management}). The following figure \ref{fig:twai-error-state-transition} shows the relation between the error states, values and bits, and error state related interrupts.

\begin{figure}[ht!]
    \centering
    \fbox{\includegraphics[width=0.8\textwidth]{\modulefiles/figures/36-TWAI-error-state-transition__EN}}
    \caption{Error State Transition}
    \label{fig:twai-error-state-transition}
\end{figure}


\subsubsection{Error Warning Limit}

The Error Warning Limit (EWL) is a configurable threshold value for the TEC and REC, which will trigger an interrupt when exceeded. The EWL is intended to serve as a warning about severe TWAI bus errors, and is triggered before the TWAI controller enters the Error Passive state. The EWL is configured in \\\hyperref[regdesc:TWAIERRWARNINGLIMITREG]{TWAI\_ERR\_WARNING\_LIMIT\_REG} and can only be configured whilst the TWAI controller is in Reset Mode. The \hyperref[regdesc:TWAIERRWARNINGLIMITREG]{TWAI\_ERR\_WARNING\_LIMIT\_REG} has a default value of 96.

When the values of TEC and/or REC are larger than or equal to the EWL value, the \hyperref[fielddesc:TWAIERRST]{TWAI\_ERR\_ST} bit is immediately set to 1. Likewise, when the values of both the TEC and REC are smaller than the EWL value, the \hyperref[fielddesc:TWAIERRST]{TWAI\_ERR\_ST} bit is immediately reset to 0. The Error Warning Interrupt is triggered whenever the value of the \hyperref[fielddesc:TWAIERRST]{TWAI\_ERR\_ST} bit (or the \hyperref[fielddesc:TWAIBUSOFFST]{TWAI\_BUS\_OFF\_ST}) changes.


\subsubsection{Error Passive}

The TWAI controller is in the Error Passive state when the TEC or REC value exceeds 127. Likewise, when both the TEC and REC are less than or equal to 127, the TWAI controller enters the Error Active state. The Error Passive Interrupt is triggered whenever the TWAI controller transitions from the Error Active state to the Error Passive state or vice versa.

\subsubsection{Bus-Off and Bus-Off Recovery}

The TWAI controller enters the Bus-Off state when the TEC value exceeds 255. On entering the Bus-Off state, the TWAI controller will automatically do the following:

\begin{itemize}
    \item Set REC to 0
    \item Set TEC to 127
    \item Set the \hyperref[fielddesc:TWAIBUSOFFST]{TWAI\_BUS\_OFF\_ST} bit to 1
    \item Enter Reset Mode 
\end{itemize}

The Error Warning Interrupt is triggered whenever the value of the \hyperref[fielddesc:TWAIBUSOFFST]{TWAI\_BUS\_OFF\_ST} bit (or the \hyperref[fielddesc:TWAIERRST]{TWAI\_ERR\_ST} bit) changes.

To return to the Error Active state, the TWAI controller must undergo Bus-Off Recovery. Bus-Off Recovery requires the TWAI controller to observe 128 occurrences of 11 consecutive recessive bits on the bus. To initiate Bus-Off Recovery (after entering the Bus-Off state), the TWAI controller should enter Operation Mode by setting the \hyperref[fielddesc:TWAIRESETMODE]{TWAI\_RESET\_MODE} bit to 0.
The TEC tracks the progress of Bus-Off Recovery by decrementing the TEC each time when the TWAI controller observes 11 consecutive recessive bits. When Bus-Off Recovery has completed (i.e., TEC has decremented from 127 to 0), the \hyperref[fielddesc:TWAIBUSOFFST]{TWAI\_BUS\_OFF\_ST} bit will automatically be reset to 0, thus triggering the Error Warning Interrupt.


\subsection{Error Code Capture}

The Error Code Capture (ECC) feature allows the TWAI controller to record the error type and bit position of a TWAI bus error in the form of an error code. Upon detecting a TWAI bus error, the Bus Error Interrupt is triggered and the error code is recorded in \hyperref[regdesc:TWAIERRCODECAPREG]{TWAI\_ERR\_CODE\_CAP\_REG}. Subsequent bus errors will trigger the Bus Error Interrupt, but their error codes will not be recorded until the current error code is read from the \hyperref[regdesc:TWAIERRCODECAPREG]{TWAI\_ERR\_CODE\_CAP\_REG}.

The following Table \ref{tab:twai-bit-info-error-code-cap-reg} shows the fields of the \hyperref[regdesc:TWAIERRCODECAPREG]{TWAI\_ERR\_CODE\_CAP\_REG}:

\begin{longtable}[c]{|p{1.4cm}|p{1.4cm}|p{1.4cm}|p{1.4cm}|p{1.4cm}|p{1.4cm}|p{1.4cm}|p{1.4cm}|p{1.4cm}|}
    \caption{Bit Information of \hyperref[regdesc:TWAIERRCODECAPREG]{TWAI\_ERR\_CODE\_CAP\_REG} (0x30)}
    \label{tab:twai-bit-info-error-code-cap-reg}\\
    \hline
    \rowcolor{lightgray}
    \textbf{Bit 31-8} & \textbf{Bit 7} & \textbf{Bit 6} & \textbf{Bit 5} & \textbf{Bit 4} & \textbf{Bit 3} & \textbf{Bit 2} & \textbf{Bit 1} & \textbf{Bit 0}\\\hline
    Reserved & ERRC.1$^1$ & ERRC.0$^1$ & DIR$^2$ & SEG.4$^3$ & SEG.3$^3$ & SEG.2$^3$ & SEG.1$^3$ & SEG.0$^3$\\
    \hline
\end{longtable}
\vspace{-1.5em}
{\bfseries Notes:}\\
\begin{itemize}
    \item ERRC: The Error Code (ERRC) indicates the type of bus error: 00 for bit error, 01 for format error, 10 for stuff error, and 11 for other types of error.
    \item DIR: The Direction (DIR) indicates whether the TWAI controller was transmitting or receiving when the bus error occurred: 0 for transmitter, 1 for receiver.
    \item SEG: The Error Segment (SEG) indicates the segment of the TWAI message at which the bus error occurred.
\end{itemize}

The following Table \ref{tab:twai-bit-info-seg4-seg0} shows how to interpret the SEG.0 to SEG.4 bits.

\begin{longtable}[c]{|p{2cm}<{\centering}|p{2cm}<{\centering}|p{2cm}<{\centering}|p{2cm}<{\centering}|p{2cm}<{\centering}|p{4cm}|}
    \caption{Bit Information of Bits SEG.4 - SEG.0}
    \label{tab:twai-bit-info-seg4-seg0}\\
    \hline
    \rowcolor{lightgray}
   \textbf{Bit SEG.4} & \textbf{Bit SEG.3} & \textbf{Bit SEG.2} & \textbf{Bit SEG.1} & \textbf{Bit SEG.0} & \textbf{Description} \\\hline
   \endfirsthead
   
   \multicolumn{6}{c}{\bfseries \tablename\ \thetable{} -- cont'd from previous page}\\\hline
   
   \rowcolor{lightgray}
   \textbf{Bit SEG.4} & \textbf{Bit SEG.3} & \textbf{Bit SEG.2} & \textbf{Bit SEG.1} & \textbf{Bit SEG.0} & \textbf{Description} \\\hline
   \endhead
   \multicolumn{6}{r}{\textbf{Cont'd on next page}} \\
   \endfoot
   \endlastfoot
   
    0 & 0 & 0 & 1 & 1 & start of frame
    \\\hline
    0 & 0 & 0 & 1 & 0 & ID.28 \verb+~+ ID.21
    \\\hline
    0 & 0 & 1 & 1 & 0 & ID.20 \verb+~+ ID.18
    \\\hline
    0 & 0 & 1 & 0 & 0 & bit SRTR
    \\\hline
    0 & 0 & 1 & 0 & 1 & bit IDE
    \\\hline
    0 & 0 & 1 & 1 & 1 & ID.17 \verb+~+ ID.13
    \\\hline
    0 & 1 & 1 & 1 & 1 & ID.12 \verb+~+ ID.5
    \\\hline
    0 & 1 & 1 & 1 & 0 & ID.4 \verb+~+ ID.0
    \\\hline
    0 & 1 & 1 & 0 & 0 & bit RTR
    \\\hline
    0 & 1 & 1 & 0 & 1 & reserved bit 1
    \\\hline
    0 & 1 & 0 & 0 & 1 & reserved bit 0
    \\\hline
    0 & 1 & 0 & 1 & 1 & data length code
    \\\hline
    0 & 1 & 0 & 1 & 0 & data field
    \\\hline
    0 & 1 & 0 & 0 & 0 & CRC sequence
    \\\hline
    1 & 1 & 0 & 0 & 0 & CRC delimiter
    \\\hline
    1 & 1 & 0 & 0 & 1 & ACK slot
    \\\hline
    1 & 1 & 0 & 1 & 1 & ACK delimiter
    \\\hline
    1 & 1 & 0 & 1 & 0 & end of frame
    \\\hline
    1 & 0 & 0 & 1 & 0 & intermission
    \\\hline
    1 & 0 & 0 & 0 & 1 & active error flag
    \\\hline
    1 & 0 & 1 & 1 & 0 & passive error flag
    \\\hline
    1 & 0 & 0 & 1 & 1 & tolerate dominant bits
    \\\hline
    1 & 0 & 1 & 1 & 1 & error delimiter
    \\\hline
    1 & 1 & 1 & 0 & 0 & overload flag
    \\\hline
\end{longtable}
\vspace{-1.5em}
{\bfseries Notes:}\\
\begin{itemize}
    \item Bit SRTR: under Standard Frame Format.
    \item Bit IDE: Identifier Extension Bit, 0 for Standard Frame Format.

\end{itemize}

\subsection{Arbitration Lost Capture}

The Arbitration Lost Capture (ALC) feature allows the TWAI controller to record the bit position where it loses arbitration. When the TWAI controller loses arbitration, the bit position is recorded in \hyperref[regdesc:TWAIARBLOSTCAPREG]{TWAI\_ARB\_LOST\_CAP\_REG} and the Arbitration Lost Interrupt is triggered.

Subsequent losses in arbitration will trigger the Arbitration Lost Interrupt, but will not be recorded in \hyperref[regdesc:TWAIARBLOSTCAPREG]{TWAI\_ARB\_LOST\_CAP\_REG} until the current Arbitration Lost Capture is read from the \hyperref[regdesc:TWAIERRCODECAPREG]{TWAI\_ERR\_CODE\_CAP\_REG}.

Table \ref{tab:twai-bit-info-arb-lost-cap-reg} illustrates bits and fields of \hyperref[regdesc:TWAIERRCODECAPREG]{TWAI\_ERR\_CODE\_CAP\_REG} whilst Figure \ref{fig:twai-arbitration-lost-position} illustrates the bit positions of a TWAI message.

\begin{figure}[ht]
    \centering
    \fbox{\includegraphics[width=0.9\textwidth]{\modulefiles/figures/36-TWAI-arbitration-lost-position__EN}}
    \caption{Positions of Arbitration Lost Bits}
    \label{fig:twai-arbitration-lost-position}
\end{figure}

\begin{longtable}[c]{|p{2.3cm}|p{2.3cm}|p{2.3cm}|p{2.3cm}|p{2.3cm}|p{2.3cm}|}
    \caption{Bit Information of \hyperref[regdesc:TWAIARBLOSTCAPREG]{TWAI\_ARB\_LOST\_CAP\_REG} (0x2c)}
    \label{tab:twai-bit-info-arb-lost-cap-reg}
    \\\hline
    \rowcolor{lightgray}
    \textbf{Bit 31-5} & \textbf{Bit 4} & \textbf{Bit 3} & \textbf{Bit 2} & \textbf{Bit 1} & \textbf{Bit 0}\\ \hline
   Reserved & BITNO.4$^1$ & BITNO.3$^1$ & BITNO.2$^1$ & BITNO.1$^1$ & BITNO.0$^1$\\\hline
\end{longtable}
\vspace{-1.5em}
{\bfseries Notes: }\\
\begin{itemize}
    \item BITNO: Bit Number (BITNO) indicates the nth bit of a TWAI message where arbitration was lost.
\end{itemize}

\subsection{Transceiver Auto-Standby}
\label{sec:twai-transceiver-auto-standby}
It is common for TWAI transceivers to support a standby mode to lower power consumption. TWAI transceivers will generally expose a standby signal that is asserted by the connected TWAI controller, thus allowing the controller to place the transceiver into standby when appropriate (e.g., when the bus will be idle for an extended period of time). Transceivers will exit standby mode if the controller de-asserts the standby signal, or if the transceiver detects bus activity (also known as a wake-up feature).

\chipname{}'s TWAI controller supports both hardware control (i.e., automatic) and software control (i.e., manual) of the standby signal to control the switching of TWAI transceivers connected to the chip. When hardware controlled, the TWAI controller will automatically assert the standby signal when the bus remains idle for longer than a configurable amount of time. When software controlled, the standby signal can be manually asserted/de-asserted directly by software.

\begin{itemize}
    \item Hardware output:
    \begin{enumerate}
        \item Set the \hyperref[fielddesc:TWAIHWSTANDBYEN]{TWAI\_HW\_STANDBY\_EN} field in the \hyperref[regdesc:TWAIHWCFGREG]{TWAI\_HW\_CFG\_REG} register to enable standby function for hardware.
        \item Configure the \hyperref[regdesc:TWAIHWSTANDBYCNTREG]{TWAI\_HW\_STANDB\_CNT\_REG} register. This register indicates the time required before hardware trigger the standby signal after entering idle status, in which the value indicates the number of cycles of the TWAI controller operating clock (40 MHz by default).
    \end{enumerate}
    \item Software output:
    \begin{enumerate}
        \item Set the \hyperref[fielddesc:TWAISWSTANDBYEN]{TWAI\_SW\_STANDBY\_EN} field in the \hyperref[regdesc:TWAISWSTANDBYCFGREG]{TWAI\_SW\_STANDBY\_CFG\_REG} register to generate standby signal in the TWAI controller.
    \end{enumerate}
\end{itemize}

The standby signal generated using either of the above methods will be pulled down (cleared) when either of the following conditions is met:
\begin{enumerate}
    \item The standby signal will be automatically cleared when the TWAI controller exits the idle status.
    \item Users can also pull down the standby signal by setting the \hyperref[fielddesc:TWAISWSTANDBYCLR]{TWAI\_SW\_STANDBY\_CLR} field in the \hyperref[regdesc:TWAISWSTANDBYCFGREG]{TWAI\_SW\_STANDBY\_CFG\_REG} register.
\end{enumerate}


%%%%%%%%%%%%%%%%%%%%%%%%
%%% Register Summary %%%
%%%%%%%%%%%%%%%%%%%%%%%%

% >>> Update the sample text below and insert
%     the info on register summary into the subfile <<<

\clearpage
\hypertarget{twai-reg-summ}{}
\section{Register Summary}
\label{sec:twai-reg-summ}

'|' here means separate line to distinguish between TWAI working modes discussed in Section \ref{sec:twai-modes} \textit{\nameref{sec:twai-modes}}. The left describes the access in Operation Mode. The right belongs to Reset Mode and is marked in red. The addresses in this section are relative to Two-wire Automotive Interface base address provided in Table \ref{tab:sysmem-base-address} in Chapter \ref{mod:sysmem} \textit{\nameref{mod:sysmem}}.

The abbreviations given in Column \textbf{Access} are explained in Section \hyperref[glossary-access-types]{\textit{Access Types for Registers}}.

\subfile{\modulefiles/36-TWAI-reg-summ__EN}


%%%%%%%%%%%%%%%%%
%%% Registers %%%
%%%%%%%%%%%%%%%%%

% >>> Update the sample text below and insert
%      the info on registers into the subfile <<<

\clearpage
\section{Registers}

'|' here means separate line. The left describes the access in Operation Mode. The right belongs to Reset Mode with red color. The addresses in this section are relative to Two-wire Automotive Interface base address (each TWAI 0 and TWAI 1 has an individual base address) provided in Table \ref{tab:sysmem-base-address} in Chapter \ref{mod:sysmem} \textit{\nameref{mod:sysmem}}.

\subfile{\modulefiles/36-TWAI-reg__EN}


\end{document}
